%% Concurrency System
\chapter{Concurrency System}
\label{ch:concurrency}

Lambdust provides lightweight, fault-tolerant concurrency through an actor model with message passing, futures for asynchronous computation, and software transactional memory for shared state management.

\section{Concurrency Overview}
\label{sec:concurrency-overview}

The concurrency system is built on several key abstractions:

\begin{itemize}
\item \textbf{Actors}: Lightweight concurrent entities that communicate via messages
\item \textbf{Futures}: Asynchronous computations with eventual values
\item \textbf{Channels}: Typed communication channels for CSP-style concurrency  
\item \textbf{STM}: Software Transactional Memory for composable shared state
\item \textbf{Parallelism}: Data parallel operations and work-stealing
\end{itemize}

\section{Actor Model}
\label{sec:actors}

\indexentry{actor}
\indexentry{spawn}
\indexentry{send}
\indexentry{receive}

Actors are the fundamental unit of concurrency in Lambdust:

\section{Actor Creation}
\label{sec:actor-creation}

\begin{syntaxspec}{(spawn computation)}
Creates a new actor running \var{computation} and returns its process identifier (PID).
\end{syntaxspec}

\begin{example}[Actor Creation]
\begin{lambdustcode}
;; Simple counter actor
(define (counter-actor initial-count)
  (letrec ((loop (lambda (count)
                   (receive  
                     (('get sender)
                      (send sender count)
                      (loop count))
                     (('increment sender)
                      (send sender (+ count 1)) 
                      (loop (+ count 1)))
                     (('set value sender)
                      (send sender 'ok)
                      (loop value))))))
    (loop initial-count)))

;; Spawn the actor
(define counter-pid (spawn (counter-actor 0)))

;; Echo actor that reflects messages
(define echo-pid 
  (spawn 
    (letrec ((loop (lambda ()
                     (receive
                       ((msg sender)
                        (send sender msg) 
                        (loop))))))
      (loop))))
\end{lambdustcode}
\end{example}

\section{Message Passing}
\label{sec:message-passing}

\begin{syntaxspec}{(send pid message)}
Sends \var{message} asynchronously to the actor identified by \var{pid}.
\end{syntaxspec}

\begin{syntaxspec}{(receive \arbno{(pattern body)})}
Receives messages matching patterns, blocking until a match is found.
\end{syntaxspec}

\begin{example}[Message Passing]
\begin{lambdustcode}
;; Send messages to counter actor
(send counter-pid `(get ,self))
(send counter-pid `(increment ,self))
(send counter-pid `(set 42 ,self))

;; Receive responses  
(receive
  ((count) (display "Current count: ") (display count) (newline)))

;; Pattern matching in receive
(define (server-loop)
  (receive
    ;; Handle different request types
    (('add x y client)
     (send client (+ x y))
     (server-loop))
    
    (('multiply x y client) 
     (send client (* x y))
     (server-loop))
    
    (('ping client)
     (send client 'pong)
     (server-loop))
    
    (('shutdown)
     (display "Server shutting down") (newline))
    
    ;; Catch-all for unknown messages
    ((unknown-msg)
     (display "Unknown message: ") (display unknown-msg) (newline)
     (server-loop))))

;; Request-response pattern
(define (call pid request timeout)
  (send pid `(,request ,self))
  (receive
    ((response) response))
  timeout)    ; Handle timeout
\end{lambdustcode}
\end{example}

\section{Actor Hierarchies and Supervision}
\label{sec:supervision}

\indexentry{supervision}
\indexentry{fault tolerance}

Actors can supervise other actors for fault tolerance:

\begin{example}[Actor Supervision]
\begin{lambdustcode}
;; Supervisor that manages worker actors
(define (supervisor-actor worker-count worker-fn)
  (let ((workers (map (lambda (i) (spawn (worker-fn i))) 
                      (range 0 worker-count))))
    (letrec ((loop (lambda (active-workers)
                     (receive
                       ;; Worker crashed, restart it
                       (('EXIT pid reason)
                        (display "Worker crashed: ") (display reason) (newline)
                        (let ((new-worker (spawn (worker-fn (length active-workers)))))
                          (loop (cons new-worker (remove pid active-workers)))))
                       
                       ;; Distribute work to workers  
                       (('work task client)
                        (if (null? active-workers)
                            (send client 'no-workers-available)
                            (begin 
                              (send (car active-workers) `(,task ,client))
                              (loop (append (cdr active-workers) (list (car active-workers)))))))
                       
                       ;; Shutdown all workers
                       (('shutdown)
                        (for-each (lambda (worker) (send worker 'shutdown)) active-workers)
                        'supervisor-shutdown)))))
      (loop workers))))

;; Worker actor that can fail
(define (worker-actor id)
  (letrec ((loop (lambda ()
                   (receive  
                     (('heavy-computation x client)
                      ;; Simulate potential failure
                      (if (= (random 10) 0)
                          (error "Worker failure") 
                          (send client (* x x)))
                      (loop))
                     (('shutdown) 
                      'worker-shutdown)))))
    (loop)))

;; Create supervised worker pool
(define supervisor (spawn (supervisor-actor 3 worker-actor)))
\end{lambdustcode}
\end{example}

\section{Futures and Promises}
\label{sec:futures}

\indexentry{future}
\indexentry{promise}
\indexentry{async}
\indexentry{await}

Futures provide asynchronous computation with eventual values:

\section{Future Creation}
\label{sec:future-creation}

\begin{syntaxspec}{(future computation)}
Creates a future that will eventually contain the result of \var{computation}.
\end{syntaxspec}

\begin{syntaxspec}{(promise)}
Creates a promise/future pair for manual completion.
\end{syntaxspec}

\begin{example}[Futures]
\begin{lambdustcode}
;; Asynchronous computation
(define future-result 
  (future 
    (begin 
      (sleep 1000)    ; Simulate long computation
      (+ 20 22))))

;; Continue with other work
(display "Doing other work...") (newline)

;; Wait for result when needed
(define result (await future-result))
(display "Result: ") (display result) (newline)

;; Multiple futures
(define f1 (future (expensive-computation-1)))
(define f2 (future (expensive-computation-2)))  
(define f3 (future (expensive-computation-3)))

;; Wait for all to complete
(define results (map await (list f1 f2 f3)))

;; Promise for manual completion
(define-values (promise resolver) (make-promise))

;; In another thread/actor
(spawn 
  (lambda () 
    (let ((data (fetch-from-network)))
      (resolve! resolver data))))

;; Wait for the promise
(define network-data (await promise))
\end{lambdustcode}
\end{example}

\section{Future Composition}
\label{sec:future-composition}

\begin{example}[Future Composition]
\begin{lambdustcode}
;; Future chaining with then
(define chained-future
  (future-then 
    (future (read-file "input.txt"))
    (lambda (content) 
      (future (process-content content)))
    (lambda (processed)
      (future (write-file "output.txt" processed)))))

;; Future combination  
(define combined-future
  (future-combine
    (list (future (computation-1))
          (future (computation-2))
          (future (computation-3)))
    (lambda (results) 
      (apply + results))))

;; Racing futures (first to complete wins)
(define fastest-result
  (future-race 
    (list (future (slow-network-request))
          (future (cache-lookup))
          (future (backup-computation)))))

;; Timeout handling
(define result-with-timeout
  (future-timeout 
    (future (potentially-slow-operation))
    5000    ; 5 second timeout
    'timeout-value))

;; Error handling in futures
(define safe-future
  (future-catch
    (future (risky-computation))
    (lambda (error)
      (begin 
        (log-error error)
        'default-value))))
\end{lambdustcode}
\end{example}

\section{Channels}
\label{sec:channels}

\indexentry{channel}
\indexentry{CSP}

Channels provide CSP-style communication:

\begin{example}[Channels]
\begin{lambdustcode}
;; Create typed channels
(define numbers-channel :: (Channel Natural))
(define numbers-channel (make-channel))

;; Producer process
(spawn 
  (lambda ()
    (do ((i 0 (+ i 1)))
        ((>= i 10))
      (channel-put! numbers-channel i)
      (sleep 100))))

;; Consumer process  
(spawn
  (lambda () 
    (let loop ()
      (let ((value (channel-get numbers-channel)))
        (display "Received: ") (display value) (newline)
        (loop)))))

;; Buffered channels
(define buffered-channel (make-channel 100))  ; Buffer size 100

;; Select on multiple channels  
(define (multi-channel-select ch1 ch2 ch3)
  (select  
    ((ch1 value) 
     (begin (display "From ch1: ") (display value) (newline)))
    ((ch2 value)
     (begin (display "From ch2: ") (display value) (newline))) 
    ((ch3 value)
     (begin (display "From ch3: ") (display value) (newline)))
    (timeout 1000
     (display "Timeout occurred") (newline))))

;; Channel transformations
(define (channel-map f input-ch)
  (let ((output-ch (make-channel)))
    (spawn 
      (lambda ()
        (let loop ()
          (let ((value (channel-get input-ch)))
            (channel-put! output-ch (f value))
            (loop)))))
    output-ch))

;; Pipeline processing
(define (pipeline input-ch . transforms)
  (fold-left channel-map input-ch transforms))
\end{lambdustcode}
\end{example}

\section{Software Transactional Memory}
\label{sec:stm}

\indexentry{STM}
\indexentry{transaction}
\indexentry{atomic}

STM provides composable shared state management:

\section{Transactional Variables}
\label{sec:tvars}

\begin{example}[Transactional Variables]
\begin{lambdustcode}
;; Create transactional variables
(define account1 (make-tvar 1000))    ; Account with $1000
(define account2 (make-tvar 500))     ; Account with $500

;; Atomic transaction  
(define (transfer amount from-account to-account)
  (atomic
    (let ((from-balance (tvar-read from-account))
          (to-balance (tvar-read to-account)))
      (if (>= from-balance amount)
          (begin 
            (tvar-write! from-account (- from-balance amount))
            (tvar-write! to-account (+ to-balance amount))
            #t)    ; Success
          #f))))   ; Insufficient funds

;; Transfer money atomically
(if (transfer 200 account1 account2)
    (display "Transfer successful")
    (display "Transfer failed"))

;; Compose transactions
(define (multi-transfer transfers)
  (atomic 
    (and-map (lambda (transfer-spec)
               (apply transfer transfer-spec))
             transfers)))

;; Retry on failure
(define (wait-for-funds account minimum)
  (atomic
    (let ((balance (tvar-read account)))
      (if (>= balance minimum)
          balance
          (retry)))))    ; Retry when condition met
\end{lambdustcode}
\end{example}

\section{Transaction Composition}
\label{sec:transaction-composition}

\begin{example}[Transaction Composition]
\begin{lambdustcode}
;; Composable transactions
(define (deposit amount account)
  (atomic 
    (let ((balance (tvar-read account)))
      (tvar-write! account (+ balance amount)))))

(define (withdraw amount account)
  (atomic
    (let ((balance (tvar-read account)))
      (if (>= balance amount)
          (begin 
            (tvar-write! account (- balance amount))
            amount)
          (error "Insufficient funds")))))

;; Compose into larger transaction
(define (compound-transaction)
  (atomic
    ;; All or nothing - if any part fails, entire transaction aborts
    (deposit 100 account1)
    (withdraw 50 account1)  
    (deposit 50 account2)
    (let ((total (+ (tvar-read account1) (tvar-read account2))))
      (if (< total 1000)
          (error "Total too low")
          total))))

;; Alternative transactions (OrElse)  
(define (try-withdraw-from-either amount acc1 acc2)
  (or-else 
    (withdraw amount acc1)     ; Try first account
    (withdraw amount acc2)))   ; Fallback to second account

;; Conditional transactions
(define (conditional-transfer condition from to amount)
  (atomic
    (if (condition)
        (transfer amount from to)
        (retry))))    ; Wait for condition
\end{lambdustcode}
\end{example}

\section{Parallel Processing}
\label{sec:parallel}

\indexentry{parallel}
\indexentry{work stealing}

Lambdust provides data parallelism and work-stealing:

\begin{example}[Parallel Processing]
\begin{lambdustcode}
;; Parallel map
(define (parallel-map f lst)
  (let ((chunks (partition lst (processor-count))))
    (let ((futures (map (lambda (chunk) (future (map f chunk))) chunks)))
      (apply append (map await futures)))))

;; Parallel fold  
(define (parallel-fold f identity lst)
  (if (<= (length lst) 100)   ; Sequential threshold
      (fold-left f identity lst)
      (let ((mid (quotient (length lst) 2)))
        (let ((left (take lst mid))
              (right (drop lst mid)))
          (let ((left-future (future (parallel-fold f identity left)))
                (right-future (future (parallel-fold f identity right))))
            (f (await left-future) (await right-future)))))))

;; Work-stealing pool
(define (work-stealing-map f lst)
  (let ((work-queue (make-work-queue lst))
        (result-queue (make-channel))
        (worker-count (processor-count)))
    
    ;; Spawn worker threads
    (do ((i 0 (+ i 1)))
        ((>= i worker-count))
      (spawn 
        (lambda ()
          (let loop ()
            (let ((work (work-queue-steal! work-queue)))
              (if work
                  (begin 
                    (channel-put! result-queue (f work))
                    (loop))
                  'worker-done))))))
    
    ;; Collect results
    (let ((results '()))
      (do ((i 0 (+ i 1)))
          ((>= i (length lst)) (reverse results))
        (set! results (cons (channel-get result-queue) results))))))

;; GPU parallel computation (when available)
(define (gpu-parallel-map f lst)
  (if (gpu-available?)
      (gpu-compute (gpu-compile-function f) (list->gpu-array lst))
      (parallel-map f lst)))   ; Fallback to CPU parallelism
\end{lambdustcode}
\end{example}

\section{Distributed Computing}
\label{sec:distributed}

\indexentry{distributed}
\indexentry{location transparency}

Actors support location transparency for distributed computing:

\begin{example}[Distributed Computing]
\begin{lambdustcode}
;; Remote actor spawning
(define remote-node "worker-node-1.cluster.local")
(define remote-actor (spawn-remote remote-node (worker-computation)))

;; Location-transparent message passing
(send remote-actor `(process-data ,large-dataset ,self))

;; Distributed supervision  
(define (distributed-supervisor nodes worker-fn)
  (let ((workers (map (lambda (node) 
                        (spawn-remote node (worker-fn))) 
                      nodes)))
    (letrec ((supervise (lambda (active-workers)
                          (receive
                            ;; Handle node failures
                            (('NODE-DOWN node)
                             (display "Node failed: ") (display node) (newline)
                             (let ((replacement (spawn-remote (backup-node) (worker-fn))))
                               (supervise (cons replacement 
                                               (filter (lambda (w) (not (on-node? w node)))
                                                      active-workers)))))
                            
                            ;; Distribute work across cluster
                            (('work task client)
                             (let ((worker (least-loaded active-workers)))
                               (send worker `(,task ,client)))
                             (supervise active-workers))))))
      (supervise workers))))

;; Fault-tolerant distributed computation
(define (distributed-map-reduce map-fn reduce-fn data)
  (let ((nodes (available-nodes)))
    ;; Distribute data across nodes
    (let ((chunks (distribute-data data nodes)))
      ;; Spawn mappers on each node
      (let ((map-futures 
              (map (lambda (node-chunk)
                     (let ((node (car node-chunk))
                           (chunk (cdr node-chunk)))
                       (future-remote node (map map-fn chunk))))
                   chunks)))
        ;; Collect and reduce results
        (let ((map-results (map await map-futures)))
          (fold-left reduce-fn (car map-results) (cdr map-results)))))))

;; Network partition handling
(define (partition-tolerant-operation data)
  (with-network-timeout 5000
    (lambda ()
      (distributed-computation data))
    (lambda ()
      ;; Fallback to local computation during partition
      (local-computation data))))
\end{lambdustcode}
\end{example}

\section{Concurrency Types}
\label{sec:concurrency-types}

\indexentry{concurrency types}
\indexentry{actor type}

The type system tracks concurrency-related properties:

\begin{example}[Concurrency Types]
\begin{lambdustcode}
;; Actor types
(define counter-actor :: (Actor (U ('get ActorId) ('increment ActorId) ('set Natural ActorId))))

;; Future types  
(define computation :: (Future Natural))
(define async-string :: (Future String))

;; Channel types
(define numbers :: (Channel Natural))
(define messages :: (Channel String)) 

;; STM types
(define account :: (TVar Natural))
(define shared-counter :: (TVar Natural))

;; Effect tracking for concurrency
(define concurrent-computation :: (-> Unit ~> {Async, STM} Natural))
(define (concurrent-computation)
  (let ((counter (make-tvar 0)))
    (let ((f1 (future (atomic (tvar-modify! counter (lambda (x) (+ x 1))))))
          (f2 (future (atomic (tvar-modify! counter (lambda (x) (+ x 1)))))))
      (await f1)
      (await f2)
      (atomic (tvar-read counter)))))

;; Actor behavior types
(define (typed-server :: (-> (Channel (Request Natural String)) Unit ~> {Actor} Unit))
  (lambda (request-channel)
    (let loop ()
      (let ((request (channel-get request-channel)))
        (send (request-client request) (process-request request))
        (loop)))))

;; Deadlock detection in types
(define (potential-deadlock :: (-> (TVar Natural) (TVar Natural) ~> {STM, Deadlock} Unit))
  (lambda (tvar1 tvar2)
    ;; Type system can warn about potential circular dependencies
    (atomic 
      (let ((val1 (tvar-read tvar1)))
        (let ((val2 (tvar-read tvar2)))
          (tvar-write! tvar1 val2)
          (tvar-write! tvar2 val1))))))
\end{lambdustcode}
\end{example}

\section{Performance and Optimization}
\label{sec:concurrency-performance}

\indexentry{concurrency performance}
\indexentry{green threads}

The concurrency implementation is optimized for performance:

\begin{example}[Concurrency Performance]
\begin{lambdustcode}
;; Lightweight actors (green threads)
(define actors (map (lambda (i) (spawn (simple-actor i))) (range 0 10000)))
;; Can spawn thousands of actors with minimal overhead

;; Work-stealing scheduler
(define (cpu-intensive-work data)
  ;; Automatically distributed across available cores
  (parallel-map expensive-function data))

;; Lock-free data structures  
(define lock-free-counter (make-atomic-counter 0))
(define (increment-counter!)
  (atomic-increment! lock-free-counter))  ; No locks, CAS operations

;; Actor batching for high throughput
(define (batch-actor batch-size)
  (let ((batch '()))
    (letrec ((loop (lambda ()
                     (receive
                       ((msg)
                        (set! batch (cons msg batch))
                        (if (>= (length batch) batch-size)
                            (begin 
                              (process-batch batch)
                              (set! batch '())
                              (loop))
                            (loop)))))))
      (loop))))

;; Memory-efficient message passing
(define (zero-copy-send pid large-data)
  ;; Large data shared between actors without copying
  (send-shared pid large-data))

;; NUMA-aware scheduling  
(define (numa-spawn computation)
  ;; Spawn on local NUMA node for better memory locality
  (spawn-local computation))
\end{lambdustcode}
\end{example}

\section{Debugging and Monitoring}
\label{sec:concurrency-debugging}

\indexentry{concurrency debugging}
\indexentry{deadlock detection}

Tools for debugging concurrent programs:

\begin{example}[Concurrency Debugging]
\begin{lambdustcode}
;; Actor monitoring
(define (monitored-actor computation)
  (spawn-monitored computation
    (lambda (pid state)
      (display "Actor ") (display pid) 
      (display " state: ") (display state) (newline))))

;; Deadlock detection
(enable-deadlock-detection!)
(define (potential-deadlock)
  (atomic 
    (let ((val1 (tvar-read tvar1)))
      ;; System detects potential circular wait
      (let ((val2 (tvar-read tvar2)))   ; May warn or abort
        (do-something val1 val2)))))

;; Message tracing
(enable-message-tracing!)
(send traced-actor message)   ; Logged for debugging

;; Performance profiling
(with-concurrency-profiler
  (lambda ()
    (distributed-computation large-dataset)))
;; Produces report on actor communication, load distribution, etc.

;; Race condition detection  
(define shared-data (make-atomic-box 0))
(spawn (lambda () (atomic-update! shared-data (lambda (x) (+ x 1)))))
(spawn (lambda () (atomic-update! shared-data (lambda (x) (* x 2)))))
;; Runtime can detect data races and report them
\end{lambdustcode}
\end{example}

The concurrency system provides powerful abstractions for parallel and distributed computing while maintaining safety through the type system and runtime checks.